
\begin{figure}[!t]  
    \centering  
    \includegraphics[width=1.0\textwidth]{figures/video_dit/dit_arch.pdf}  
    \caption{Architecture of the \method.} 
    \label{fig:overview_network} % 
\end{figure}


\subsection{Model Training}
\label{sec:model_training}

In this section, we will provide a detailed introduction to the architecture design of the foundational video model, specifically for the text-to-video task, along with the details of the pre-training and post-training phases.



In Fig.~\ref{fig:overview_network}, we present the architecture of \method, which is designed based on the current mainstream DiT~\citep{dit} architecture, which typically consists of three primary components: \method-VAE, diffusion transformer, and text encoder.
%
For a given video $V\in \mathbb{R}^{(1+T) \times H \times W \times 3}$, the encoder of \method-VAE encodes it from pixel space into the latent space $x \in \mathbb{R}^{1+T/4 \times H/8 \times W/8}$, which is then fed into the subsequent diffusion transformer structure.





\subsubsection{Video Diffusion Transformer}
\label{sec:dit}

\begin{wrapfigure}{r}{0.4\linewidth}
    \vspace{-0.2em}
    \includegraphics[width=0.24\textheight]{figures/video_dit/dit_block.pdf}
    \caption{Transformer block of \method.}
    \label{fig:dit_block}
\end{wrapfigure}
%
\textbf{Diffusion transformer}. The diffusion transformer mainly consists of three components: a patchifying module, transformer blocks, and an unpatchifying module.
%
Within each block, we focus on effectively modeling spatio-temporal contextual relationships and embedding text conditions alongside time steps.
%
In the patchifying module, we use a 3D convolution with a kernel size of $(1, 2, 2)$ and apply a flattening operation to convert $x$ into a sequence of features with the shape of $(B, L, D)$, where $B$ denotes the batch size, $L=(1+T/4) \times H/16 \times W/16$ represents the sequence length, and $D$ indicates the latent dimension.
%
Specifically, as illustrated in Fig.~\ref{fig:dit_block}, we employ the cross-attention mechanism to embed the input text conditions, which can ensure the model's ability to follow instructions even under long-context modeling.
%
Additionally, we employ an MLP with a Linear layer and a SiLU~\citep{SiLU} layer to process the input time embeddings and predict six modulation parameters individually. 
%
This MLP is shared across all transformer blocks, with each block learning a distinct set of biases.
% 
Through extensive experiments, we have demonstrated that this design can reduce the parameter count by approximately $25\%$ and reveal a significant performance improvement with this approach at the same parameter scale.
%


\textbf{Text encoder}. \method's architecture uses the umT5~\citep{chung2023unimax} to encode input text.
%
Through extensive experiments, we find that umT5 has several advantages in our framework: 
%
1) It possesses strong multilingual encoding capabilities, allowing it to understand both Chinese and English effectively, as well as input visual text; 
%
2) Under the same conditions, we find that umT5 outperforms other unidirectional attention mechanism LLMs in composition;
%
3) It exhibits superior convergence, with umT5 achieving faster convergence at the same parameter scale. 
%
Based on these findings, we ultimately chose umT5 as our text embedder.



\subsubsection{Pre-training}
We leverage the flow matching framework~\citep{flow-matching, esser2024sd3}  to model a unified denoising diffusion process across both image and video domains. 
We first conduct pre-training on low-resolution images, followed by multi-stage joint optimization of images and videos. 
The image-video joint training evolves using progressively upscaled data resolutions and extended temporal durations throughout the stages.

\noindent\textbf{Training objective.}
Flow matching provides a theoretically grounded framework for learning continuous-time generative processes in diffusion models. 
It circumvents iterative velocity prediction, enabling stable training via ordinary differential equations (ODEs) while maintaining equivalence to maximum likelihood objectives.
During training, given an image or video latent $x_{1}$, a random noise $x_{0} \sim \mathcal{N}(0, I)$, and a timestep $t \in[0,1]$ sampled from a logit-normal distribution, an intermediate latent $x_{t}$ is obtained as the training input.
Following Rectified Flows (RFs)~\citep{esser2024sd3}, $x_{t}$ is defined as a linear interpolation between $x_{0}$ and $x_{1}$, \ie
\begin{equation}\label{x_t}
    x_{t}=t x_{1}+(1-t) x_{0}.
\end{equation}
The groud truth velocity $v_{t}$ is
\begin{equation}\label{v_t}
v_{t}=\frac{d x_{t}}{d t}=x_{1}-x_{0}.
\end{equation}
The model is trained to predict the velocity, thus, the loss function can be formulated as the mean squared error (MSE) between the model output and $v_{t}$,

\begin{equation}\label{rf_loss}
\mathcal{L} = \mathbb{E}_{x_0, x_1, c_{txt}, t}||u(x_t, c_{txt}, t; \theta)-v_t||^2,
\end{equation}
where $c_{txt}$ is the umT5 text embedding sequence of 512 tokens long, $\theta$ is the model weights, and $u(x_t, c_{txt}, t; \theta)$ denotes the output velocity predicted by the model.


\noindent\textbf{Image pre-training.}
Our experimental analysis reveals two critical challenges in direct joint training with high-resolution images and long-duration video sequences: 
(1)~Extended sequence lengths (typically $81$ frames for $1280\times720$ video) substantially reduce training throughput. Under fixed GPU-hour budgets, this results in insufficient data throughput, thereby impeding model convergence rates. 
(2)~Excessive GPU memory consumption forces suboptimal batch sizes, inducing training instability caused by spikes in gradient variance.
To mitigate these issues, we initialize the 14B model training through low-resolution ($256~px$) text-to-image pre-training, enforcing cross-modal semantic-textual alignment and geometric structure fidelity before progressively introducing high-res video modalities.

\noindent\textbf{Image-video joint training.}
Following large-scale $256~px$ text-to-image pre-training, we implement staged joint-training with image and video data through a resolution-progressive curriculum. 
The training protocol comprises three distinct stages differentiated by spatial resolution areas:
(1)~In the first stage, we conduct joint training with $256~px$ resolution images and 5-second video clips ($192~px$ resolution, $16$ fps).
(2)~In the second stage, we initiate spatial resolution scaling by upgrading both image and video resolutions to $480px$ while maintaining a fixed 5-second video duration.
(3)~In the final phase, we escalate spatial resolution to $720 px$ for both images and 5-second video clips.

\noindent\textbf{Pre-training configurations.}
We employ efficient training at bf16-mixed precision combined with the AdamW~\citep{adamw, adam} optimizer with a weight decay coefficient of $1e^{-3}$.
The initial learning rate is set to $1e^{-4}$, and dynamically reduced triggered by plateaus of the FID and CLIP Score metrics.

\subsubsection{Post-training}

In the post-training stage, we maintain the same model architecture and optimizer configuration from the pre-training stage, initializing the network with the pre-trained checkpoint. We conduct joint training at resolutions of $480 px$ and $720 px$ using the post-training video dataset detailed in Sec.~\ref{sec:post-training data}.